\chapter{Control de versiones}
\label{Control de versiones}

Descripción del control de versiones del proyecto: el repositorio propio que se entrega, la rama de
trabajo sobre el marco de simulación de terceros y el mecanismo que mantiene sincronizadas ambas
copias.

\section{Organización del repositorio}

Estructura de directorios del entregable y criterio de qué se versiona y qué no. En particular, por
qué los conjuntos de datos y los puntos de control de los modelos quedan fuera del repositorio.

\section{Los cambios sobre el marco de simulación}

El trabajo modifica ficheros de un proyecto de terceros. Esos cambios se conservan como parches
generados contra la base de la versión oficial, y no como copias completas, para que quede
explícito qué se ha tocado y para poder rebasarlos sobre versiones posteriores.

\begin{table*}[ht]
	\centering
	\caption{Parches aplicados sobre el marco de simulación}
	\label{tab:parches}
	\makebox[\textwidth]{
		\begin{tabular}{|l|p{10cm}|}
			\cline{1-2}
			Parche & Qué cambia \\ \hline
			[PENDIENTE] & \\ \hline
		\end{tabular}
	}
\end{table*}

\section{Historial de desarrollo}

Resumen del historial de confirmaciones, con el criterio seguido para los mensajes y para el reparto
del trabajo en incrementos verificables.
